\section{Definición del alcance}

El presente protocolo se enfoca en el diseño e implementación del frontend móvil del sistema de análisis automatizado de URLs para detección de quishing, integrado dentro de un sistema compuesto por una interfaz en React Native y un motor backend desplegado en AWS. El frontend cubre el flujo completo de la experiencia de usuario e integra el Filtro F1 como primer filtro de análisis funcional end-to-end:

\begin{itemize}
    \item \textbf{F1 --- Reputación:} consumo en tiempo real del veredicto consolidado de Google Safe Browsing, PhishTank y VirusTotal; primer filtro funcional integrado de extremo a extremo.
    \item \textbf{Pantallas de resultado para F2.1} (análisis heurístico), \textbf{F2.2} (sandbox) y \textbf{F2.3} (vista previa visual segura): la interfaz presenta el desglose y los indicadores recibidos del backend; la implementación del motor de análisis de estos filtros forma parte del módulo backend.
    \item \textbf{Funcionalidades secundarias:} historial local offline (últimos 20 escaneos), reporte de QR malicioso al backend y Story Card compartible con geocodificación inversa.
\end{itemize}

El producto final es demostrable en dispositivos físicos iOS y Android; no contempla publicación en tiendas de aplicaciones.

\section{Delimitación Funcional y Técnica}

\subsection{Funcionalidades incluidas}

\begin{table}[H]
    \centering
    \caption{Delimitación Funcional y Técnica}
    \begin{tabular}{@{}ll@{}}
        \toprule
        \textbf{Módulo} & \textbf{Componente} \\
        \midrule
        Frontend & Decodificación QR nativa en tiempo real (Apple Vision / ML Kit) \\
        Frontend & Clasificación de contenido: URL vs. no-URL \\
        Frontend & Pantalla de carga (skeleton/shimmer) y resultado simple/extendido \\
        Frontend & Historial local cifrado (últimos 20 escaneos, consultable offline) \\
        Frontend & Reporte de QR malicioso al backend \\
        Frontend & Story Card compartible con ubicación geocodificada \\
        Frontend & Integración del Filtro F1 (consumo del veredicto de reputación vía backend) \\
        \bottomrule
    \end{tabular}
    \label{tab:delimitacion}
\end{table}

El motor de análisis backend (filtros F2.1, F2.2 y F2.3) forma parte del sistema integral, pero su implementación queda fuera del alcance principal del presente protocolo; el frontend consume sus resultados mediante el contrato de interfaz acordado.

\subsection{Stack tecnológico fijo}

\begin{table}[H]
    \centering
    \caption{Stack tecnológico fijo}
    \begin{tabular}{@{}ll@{}}
        \toprule
        \textbf{Capa} & \textbf{Tecnología} \\
        \midrule
        Framework móvil & React Native 0.76+ (New Architecture) \\
        Cámara y escaneo & react-native-vision-camera v4 \\
        Almacenamiento local & react-native-mmkv \\
        Navegación & React Navigation \\
        Geocodificación inversa & Nominatim / Google Maps API \\
        Reputación (F1) & Google Safe Browsing v5, PhishTank, VirusTotal v3 \\
        \bottomrule
    \end{tabular}
    \label{tab:stack}
\end{table}

La unidad de análisis es exclusivamente la URL textual extraída del QR (RFC 3986). No se analiza el símbolo QR en sí ni ningún otro tipo de contenido codificado.

\section{Restricciones}

\begin{itemize}
    \item \textbf{Temporal:} diecisiete semanas de desarrollo organizadas en seis fases, contemplando un margen adicional dentro del semestre académico para atender posibles ajustes o imprevistos.
    \item \textbf{Equipo:} dos integrantes, cada una responsable de un módulo. Sin recursos adicionales.
    \item \textbf{Económica:} el sistema opera dentro de los planes gratuitos de todas las APIs externas y del AWS Free Tier. No se contempla inversión en planes de pago.
    \item \textbf{Hardware:} la compilación para iOS requiere acceso periódico al laboratorio de macOS de la universidad.
    \item \textbf{Conectividad:} el análisis requiere conexión a internet. Sin conexión, el sistema no produce veredicto y lo comunica al usuario; no se implementa análisis offline.
\end{itemize}

\section{Exclusiones explícitas}

\begin{itemize}
    \item \textbf{Funcionales:} cuentas de usuario y autenticación; historial ilimitado; GIF animado del comportamiento del sitio (F2.3); notificaciones push; escaneo en segundo plano.
    \item \textbf{De contenido:} análisis de seguridad sobre datos WiFi, vCard o esquemas de pago codificados en QR; URLs con esquemas distintos de \texttt{http://} y \texttt{https://}.
    \item \textbf{Técnicas:} implementación del motor de análisis backend (filtros F2.1, F2.2 y F2.3); modelos de aprendizaje automático propios; análisis de HTML o JavaScript en el dispositivo; despliegue en tiendas de aplicaciones; pruebas de carga a escala.
    \item \textbf{De investigación:} estudios de usabilidad con muestras estadísticamente representativas; estudios de campo en entornos reales; análisis comparativo con soluciones comerciales.
\end{itemize}